\chapter{The Core Object Model}
\label{ch:pebble-objects}

A graph editor must distinguish changing an instruction from changing its identity.
If deleting value 4 renumbers every later value, a diagnostic that refers to value 7 may silently begin naming a different computation.
PebbleIR starts with stable slots and makes the cost of retaining them explicit.

\section{Function-local identity}
\label{sec:pebble-identity}
A \term{value ID} is an integer from 1 through the function's allocated slot count.
Zero means no value.
A block ID runs from zero through the block count minus one; block zero is entry.
An ID is meaningful only with its owning function.
The pair consisting of function identity and ID is the reference, even when the C interface passes the two pieces separately.

The \code{Pir} object contains value records, terminators, counts, and a revision.
A value record names its opcode, result type, owning block, two operands, immediate, and two predecessor fields.
Unused fields must be zero in a verified function.
This canonical form prevents two encodings of the same instruction from accumulating different stale data.
Table~\ref{tab:pebble-model} relates the containers to their jobs.

\begin{table}[htbp]
\caption{The executable object model has one ownership boundary: the function.}
\label{tab:pebble-model}
\centering
\begin{tabular}{ll}
\toprule
Object & Responsibility \\
\midrule
\code{Pir} & Own slots, blocks, and revision \\
\code{PirValue} & Describe one result or a tombstone \\
\code{PirTerm} & End one block and identify successor blocks \\
\code{PirFacts} & Record derived predecessor and dominator sets \\
\code{PirVm} & Own a separate lowered instruction array \\
\bottomrule
\end{tabular}
\end{table}

There is no module implementation yet because there are no symbols or calls to resolve.
A module extension should own a table of functions with stable function IDs, a symbol-to-function map, and explicit declaration signatures.
Adding only a vector of function pointers would leave cross-function references and destruction order unspecified.
Keep function-local IDs local when introducing that outer layer.

\section{Order without pointers}
\label{sec:pebble-order}
Within a block, live records execute in increasing ID order.
This allows construction to interleave records from different blocks, but ordinary same-block operands must have smaller IDs than their uses.
Phis are special: their operands are used on incoming edges, so a loop phi can refer to a later-allocated latch value.
The verifier requires phis to precede ordinary instructions within their block.

Algorithm~\ref{alg:pebble-enumerate} is the actual block-enumeration strategy.
It is intentionally simple.
For $B$ blocks and $N$ allocated slots, enumerating the whole function this way costs $O(BN)$ rather than $O(N)$.
At eight blocks and 128 slots, that cost is acceptable for an inspectable specimen.
An extensible implementation can replace it with per-block ID sequences while preserving the identity contract.

\begin{algorithm}[htbp]
\caption{Enumerating one block without moving value records.}
\label{alg:pebble-enumerate}
\begin{pseudocode}
\Require function $F$ and valid block ID $b$
\Ensure live values of $b$ in execution order
\For{$i\gets1$ to $F.count$}
\If{$F.values[i]$ is live and its owner block is $b$}
\State yield $i$
\EndIf
\EndFor
\State yield the separate terminator of $b$
\end{pseudocode}
\end{algorithm}

\section{A worked deletion}
\label{sec:pebble-tombstone}
The running expression allocates IDs 1 through 7.
ID 6 adds constants at IDs 4 and 5.
After folding, ID 6 becomes constant 8 and the now-unused IDs 4 and 5 become \term{tombstones}.
IDs 6 and 7 remain 6 and 7.
The count remains seven because it measures allocated identity slots, not live instructions.
The printer preserves the tombstones so reparsing preserves the same IDs.

This choice makes deletion cheap but eventually exhausts the slot limit in an edit-heavy client.
Compaction is a separate transformation: build an old-to-new map, rewrite every operand and terminator, verify the new object, and retire all old handles.
Silently compacting during deletion would violate the promise that made stable IDs useful.

\section{Exercises}
\label{sec:pebble-three-exercises}
\begin{enumerate}
\item \label{ex:pebble-id-owner} Two functions both contain value 3. Explain why a numeric ID alone cannot detect a mistaken cross-function operand.
\item \label{ex:pebble-compact} Compact live IDs $\{1,2,3,6,7\}$ into consecutive IDs. Rewrite operands of old value 7, which multiplies old values 3 and 6.
\end{enumerate}

\section{Further reading}
\label{sec:pebble-three-reading}
Morgan discusses representations in the context of transformations that repeatedly change them; read with ownership and identity obligations in mind~\cite{morgan-1998-optimizing}.
